% Predicting a Coding Agent's Next Action Is Mostly a Labeling Problem
% Workshop submission (non-archival, 4 pages + references).
%
% This preamble is self-contained so it compiles anywhere (Overleaf, TeX Live)
% without a venue style file. To switch to the official workshop template,
% replace the geometry/title block with \usepackage{neurips_2025} (or whichever
% style the workshop ships) and keep the body unchanged.
%
% !! BEFORE SUBMITTING: verify every entry in the bibliography at the bottom.
%    Those three references were carried over from the project's own literature
%    notes and have NOT been re-checked against the arXiv listings.
\documentclass[10pt]{article}
\usepackage[letterpaper,margin=1in]{geometry}
\usepackage{times}
\usepackage{booktabs}
\usepackage{amsmath}
\usepackage{amssymb}
\usepackage{url}
\usepackage[hidelinks]{hyperref}
\usepackage{caption}
\captionsetup{font=small,labelfont=bf,skip=4pt}
\usepackage{titlesec}
\titlespacing*{\section}{0pt}{1.0ex plus .2ex}{0.6ex}
\setlength{\parskip}{0pt}
\setlength{\textfloatsep}{8pt}

\title{\vspace{-2em}Predicting a Coding Agent's Next Action Is Mostly a Labeling Problem:\\
Evidence from 8{,}404 Real Developer--Agent Turns}

\author{%
  Medha Bhattacharya\\
  \texttt{medhab2005@gmail.com}
}
\date{}

\begin{document}
\maketitle
\vspace{-2em}

\begin{abstract}
\noindent
We mine 8{,}404 developer--agent turns from 251 public GitHub repositories that
commit SpecStory IDE session transcripts, label each turn with what the agent
\emph{did} (\texttt{implement}, \texttt{debug}, \texttt{explain},
\texttt{review}, \texttt{tool\_only}, \texttt{other}), and ask whether the
developer's prompt plus the conversation so far predicts the agent's next action.
Our answer is a qualified no, and the more useful finding sits one level up. Two
defensible operationalizations of the \emph{same} six-category taxonomy over the
\emph{same} turns --- a metadata-informed rule, and an LLM judge reading the full
prompt/response pair --- agree on only 42.5\% of turns ($\kappa = 0.22$), and
disagree so systematically that the corpus describes a different phenomenon
depending on which you believe (\texttt{debug}: 3.5\% of turns under the rule,
26.6\% under the LLM). Worse, sharpening two category \emph{definitions} and
re-running the same LLM judge flipped 37.5\% of its labels ($\kappa = 0.54$
against its own earlier output). Human adjudication of 98 turns favors the LLM
labels ($\kappa = 0.51$ vs.\ $0.16$) without making them authoritative. Against
those labels, every representation we test --- TF-IDF, sentence embeddings,
prompt metadata, prior-turn features, repository metadata, raw prior-turn text
--- lands in a narrow 0.34--0.38 macro-F1 band under repository-held-out
evaluation. Prior-turn context helps consistently but slightly ($+0.023$
macro-F1 across 5 paired folds, $p = 0.0005$), and nearly all of that gain comes
from one feature: the previous turn's action label. Finally, two successive
corrections to our own pipeline --- fixing a silently-truncating transcript
parser, then relabeling --- shrank that same context effect from an initially
measured $+0.175$ to $+0.023$. We release the corpus, the label sets, and the
pipeline.
\end{abstract}

\section{Introduction}

Interpreting agent behavior requires first agreeing on what an agent \emph{did}.
Much recent work on coding agents evaluates task success; less asks which
categories of action actually occur in real sessions, whether those categories
can be reliably assigned, and whether they are anticipatable from what the
developer said. This paper reports on all three, using real, unsolicited
developer--agent transcripts rather than benchmark rollouts.

We set out to build a predictor: given the current prompt and the session so
far, predict the agent's next action type --- a capability that would support
model routing, tool pre-warming, and adaptive context budgets. We found instead
that the prediction task is bounded above by a labeling problem we had treated
as a solved preprocessing step. Our contributions:

\begin{enumerate}\itemsep1pt \parskip0pt \topsep2pt
\item \textbf{A corpus and pipeline.} 8{,}404 turns across 251 repositories and
3{,}191 session logs, with session structure, repository metadata, and three
parallel label sets (rule-based, LLM-assisted, human-adjudicated sample).
\item \textbf{A label-instability result, along two axes.} Two reasonable
implementations of the same taxonomy agree on 42.5\% of turns
($\kappa = 0.22$), with disagreement concentrated in specific categories rather
than spread as noise --- and one labeler re-run after a refinement to two written
category definitions agrees with itself on only 62.5\%.
\item \textbf{A near-negative prediction result.} All representations cluster at
0.34--0.38 macro-F1 (6 classes, repository-held-out). Conversational context
yields a small but statistically detectable improvement, almost entirely
attributable to label autocorrelation; repository metadata and raw prior-turn
text add nothing.
\item \textbf{A methodological caution.} How our own headline context effect fell
from $+0.175$ to $+0.023$ macro-F1 across two corrections, and how the first of
them was caught.
\end{enumerate}

\paragraph{Related work.} Two contemporaneous studies mine SpecStory logs
similarly: Tang et al.~\cite{tang2026programming} build a behavioral taxonomy
over 11{,}579 IDE sessions, and Fang et al.~\cite{fang2026conversation}
characterize agent contributions at the project level. Both describe behavior;
neither asks whether it is predictable turn-to-turn, nor reports agreement
between alternative labelings of it. Our motivation --- route the request before
generating the response --- follows RouteLLM~\cite{ong2024routellm}, with action
type substituted for difficulty.

\section{Corpus}

\textbf{Source and extraction.} SpecStory is an IDE extension that writes
Markdown transcripts of coding-agent sessions to \texttt{.specstory/history/} in
the developer's repository. Some developers commit these logs, making real,
unprompted developer--agent conversation reachable through the GitHub Search API.
We ran six discovery queries (three repository-search signals: \texttt{specstory
in:readme}, \texttt{topic:specstory}, \texttt{path:.specstory}; three code-search
signals over the same paths), crawled each candidate for
\texttt{.specstory/history/*.md} via the Contents API, and parsed transcripts
into prompt/response turn pairs, recording prompt and response text, fenced code
blocks (count, language, length), tool-use blocks, and session identifiers.

\textbf{Cleaning.} From 68{,}943 raw turn pairs we drop empty turns, bare
acknowledgments (``ok, continue''), pasted IDE diagnostics, browser console
dumps, CI logs, slash commands, path-only messages, exact duplicates, and
extreme lengths. Raw turn counts are dominated by a few repositories --- one
contributed over 30{,}000 turns, almost certainly a looping logging artifact ---
so we cap each repository at 50 turns, downsampled stratified by action label so
a repository's behavioral \emph{mix} survives the cap. The result is 8{,}404
turns across 251 repositories and 3{,}191 session logs, with a repository Gini of
0.299 and the top 10 repositories supplying 5.9\% of turns.

\textbf{Properties.} Median prompt length is 15 words; 56.7\% of turns have
fenced code in the agent's response and 23.1\% contain tool-use blocks. Sessions
have a median of 11 turns (max 207) as recorded before the cap. 31\% of prompts
were not originally in English (predominantly Chinese); these were
machine-translated for modeling, with the original text retained. Repositories
skew small and young (median 0 stars, median age 14.5 months; TypeScript 82,
Python 50, JavaScript 28), and the corpus captures only developers who both used
SpecStory \emph{and} committed the logs publicly --- not a random sample of
AI-assisted development.

\section{Three Labelings of the Same Turns}

\texttt{action\_type} is not in the raw data; it must be constructed. We built it
three ways. \textbf{(1) Rule:} a deterministic priority-ordered rule over prompt
and response --- tool use without surfaced code $\rightarrow$ \texttt{tool\_only};
multi-block or long code $\rightarrow$ \texttt{implement}; otherwise keyword
regexes over the prompt disambiguate single-block and no-code responses.
\textbf{(2) LLM judge:} \texttt{gemini-3.1-flash-lite} reads the prompt (up to 4K
chars), the response (up to 20K chars) and response metadata, and returns a
category, a confidence and a one-sentence rationale, with in-corpus few-shot
exemplars. All 8{,}404 turns were labeled (8{,}398 successfully; mean
self-reported confidence 0.95). \textbf{(3) Human adjudication:} a stratified
212-turn sample, of which 98 turns --- those where the rule and an earlier LLM
pass disagreed --- were adjudicated by the author.

These were built in that order, and the order matters: adjudicating the 98 turns
exposed two category definitions that were doing no work, so \texttt{debug} was
redefined around ``a concrete root cause was identified'' rather than ``a stack
trace was present,'' \texttt{review} around ``no reported problem is driving
this,'' and the whole corpus was relabeled. Everything below, and every model in
Section 4, uses that final relabeling.

\begin{table}[t]
\centering
\small
\begin{tabular}{lrrrrrr}
\toprule
& \texttt{implement} & \texttt{debug} & \texttt{explain} & \texttt{other} & \texttt{review} & \texttt{tool\_only} \\
\midrule
Rule label   & 53.0\% & 3.5\%  & 19.0\% & 3.5\% & 4.2\% & 16.9\% \\
LLM label    & 39.5\% & 26.6\% & 13.5\% & 7.7\% & 6.8\% & 5.8\%  \\
\midrule
Rule recall vs.\ LLM & 0.68 & 0.10 & 0.43 & 0.44 & 0.09 & 0.55 \\
\bottomrule
\end{tabular}
\caption{Class distribution under two labelings of the same 8{,}398 turns, and
the fraction of each LLM-labeled class that the rule also assigns to that class.
Overall raw agreement 42.5\%; Cohen's $\kappa = 0.223$.}
\label{tab:labels}
\end{table}

\textbf{The labelings disagree, and not randomly.} Agreement is 42.5\%
($\kappa = 0.223$, ``fair''), and the disagreement is structural
(Table~\ref{tab:labels}): 1{,}448 turns the rule calls \texttt{implement} the LLM
calls \texttt{debug}, and 646 turns the rule calls \texttt{tool\_only} the LLM
calls \texttt{implement}. The rule recovers only 10\% of LLM-\texttt{debug} and
9\% of LLM-\texttt{review} turns. A reader given the rule labels would conclude
that debugging is a rare corner of agent work (3.5\% of turns); given the LLM
labels, that it is the second most common thing agents do (26.6\%). Both
conclusions come from the same 8{,}398 transcripts.

\textbf{A labeler also disagrees with itself when the definitions move.} We can
isolate the effect of the definition refinement above, because 299 turns were
labeled by the same model both before and after it. Only the written definitions
of \texttt{debug} and \texttt{review} changed; the model, the corpus, the
metadata and the six category names did not. 37.5\% of the labels flipped
($\kappa = 0.540$ against its own earlier output). On this sample \texttt{debug}
rose from 13.0\% to 29.8\% of turns and \texttt{tool\_only} fell from 15.7\% to
6.0\%.\footnote{This sample is class-stratified, so none of these figures are
corpus-representative; they describe the sample. The full corpus was labeled
only under the refined definitions.} The category names were fixed throughout;
only the guidance text behind them moved --- and with it, the answer to ``how
much of coding-agent work is debugging?''

\textbf{Human adjudication favors the LLM without settling the question.} On the
98 adjudicated turns, the final LLM label matched human judgment on 61.2\%
($\kappa = 0.512$) versus 31.6\% for the rule ($\kappa = 0.163$). The
pre-refinement LLM labels scored $\kappa = 0.482$ on the same turns, so the
refinement did move the labels toward the human. If we additionally assume that
the 76 unadjudicated turns where the rule and the final LLM label agree are
correct --- an optimistic assumption --- these rise to $\kappa = 0.727$ (LLM) and
$0.532$ (rule) over $n = 174$. We report both, because the gap measures how much
of our ``validated'' label quality rests on an untested assumption. Two caveats:
the 98 turns were selected as disagreements under the \emph{earlier} labeling
(27 of them now agree), so they are no longer a clean disagreement sample; and
the annotator saw the earlier LLM label while adjudicating. Either way, the best
available label set agrees with a careful human reader on roughly six of ten
hard cases. That is the ceiling the next section is scored against.

\section{Can the Next Action Be Predicted?}

\textbf{Setup.} We target the LLM labels (8{,}398 turns). Repositories --- not
turns --- are split 80/20 by \texttt{GroupShuffleSplit} (200 train / 51 test
repositories; 6{,}665 / 1{,}733 turns; zero overlap), so performance must
generalize across projects rather than memorize a project's vocabulary. No
response-side column is ever a model input, since the rule label is partly
derived from such signals: features are the current prompt plus
\emph{prior}-turn signals, computed with a strict \texttt{shift(1)} within each
session. The primary metric is macro-F1 over 6 classes. We also run 5-fold
repository-held-out \texttt{GroupKFold}, reusing the same folds across
representations so per-fold scores are paired. Models: logistic regression
(class-balanced), random forest, a 256--128 MLP, and XGBoost.

\begin{table}[t]
\centering
\small
\begin{tabular}{llcc}
\toprule
Representation & Model & Macro-F1 (5-fold CV) & Accuracy \\
\midrule
Majority class$^\dagger$ & --- & 0.086 & 0.350 \\
Prompt keyword rule$^\dagger$ & --- & 0.320 & 0.428 \\
\midrule
SBERT prompt only & logistic reg. & $0.341 \pm 0.014$ & $0.425 \pm 0.022$ \\
SBERT + prev.\ action only & logistic reg. & $0.361 \pm 0.018$ & $0.444 \pm 0.027$ \\
SBERT + metadata + history & logistic reg. & $0.364 \pm 0.017$ & $0.446 \pm 0.028$ \\
SBERT + metadata + history & MLP & $\mathbf{0.364 \pm 0.024}$ & $0.559 \pm 0.025$ \\
SBERT + metadata + history + repo & logistic reg. & $0.362 \pm 0.015$ & $0.446 \pm 0.028$ \\
SBERT + metadata + history & XGBoost & $0.308 \pm 0.009$ & $0.568 \pm 0.022$ \\
\midrule
Prompt + prior-turn \emph{text}, $N{=}1^\dagger$ & MLP & 0.314 & 0.517 \\
Prompt + prior-turn \emph{text}, $N{=}3^\dagger$ & MLP & 0.309 & 0.470 \\
\bottomrule
\end{tabular}
\caption{Repository-held-out results against the LLM labels. $\pm$ is the
standard deviation over 5 repository-held-out folds. $^\dagger$ single split
only.}
\label{tab:results}
\end{table}

\textbf{Everything lands in the same narrow band} (Table~\ref{tab:results}). The
best configuration reaches 0.364 macro-F1, prompt-only reaches 0.341, and a
hand-written keyword rule over the prompt alone reaches 0.320: a learned semantic
representation buys roughly 0.02--0.04 macro-F1 over a regex cascade. Accuracy
diverges from macro-F1 in the expected way --- XGBoost posts the best accuracy
(0.568) with the \emph{worst} macro-F1 of the linear and neural families, by
concentrating predictions on frequent classes.

\textbf{Context helps, consistently and slightly.} Adding prior-turn features to
the prompt-only model improves macro-F1 by 0.023 on every one of the 5 paired
folds (deltas $+0.030, +0.018, +0.027, +0.021, +0.020$; paired $t = 10.4$,
$p = 0.0005$; Wilcoxon $p = 0.0625$, the smallest value attainable at $n{=}5$).
The effect is real and reproducible. It is also small enough that no downstream
routing decision would change because of it.

\textbf{And the gain is mostly label autocorrelation.} Replacing the full history
bundle (previous action, whether the previous turn had code or tool use, rolling
session code rate, turn position) with the previous turn's action label
\emph{alone} recovers 0.361 of the 0.364 --- a 0.003 difference not
distinguishable at this fold count ($p = 0.10$). Sessions are sticky: the best
single predictor of what the agent does next is what it just did. That is a
weaker claim than ``conversational context is informative,'' and it is the one
the data supports.

\textbf{Repository metadata and raw prior text do not help.} Adding repository
language, stars, size, age and open-issue count changes macro-F1 by $-0.002$.
Replacing summary history features with the concatenated text of the previous $N$
turns is \emph{worse} than the prompt alone at every $N$ tested (0.314, 0.301,
0.309 for $N = 1,2,3$ vs.\ 0.357 prompt-only, same split): prior-turn text
dilutes the current prompt's embedding faster than it adds signal.

\textbf{Per-class structure.} On the best single-split model, per-class F1 is
\texttt{debug} 0.72, \texttt{implement} 0.62, \texttt{explain} 0.41,
\texttt{review} 0.26, \texttt{other} 0.23, \texttt{tool\_only} 0.04. The classes
the model predicts are the ones the labelers agree on; the ones it cannot are
exactly those Section 3 shows to be unstable. \texttt{tool\_only} is near
unpredictable from prompt text --- consistent with it being a property of the
agent's harness rather than of what the developer asked for.

\section{A Cautionary Result on Our Own Pipeline}

The first version of this pipeline reported the prompt-only $\rightarrow$
+history comparison as $0.237 \rightarrow 0.412$ macro-F1: a 74\% relative
improvement, and a publishable-looking headline. Fixing the parser took it to
$0.223 \rightarrow 0.348$ ($\Delta = +0.125$); relabeling took it to the
$+0.023$ reported above.

The \textbf{parser bug}: SpecStory splits one logical agent turn across several
consecutive Markdown blocks, and our parser read only the first, silently
emptying a large share of \texttt{agent\_response} values. The rule labeler
routed those blank responses into \texttt{other}/\texttt{review} by default,
making the affected rows trivially predictable once the previous (equally
affected) turn's label was known. It was not caught by a test; it was caught by
asking why inter-labeler agreement was anomalously near zero for exactly two
classes, sampling a disagreement, and reading the raw transcript. The
\textbf{label change} accounts for most of the remaining shrinkage, consistent
with Section 3: the rule label is partly a function of response-side surface
features that persist across adjacent turns, so \texttt{prev\_action\_type}
carried more autocorrelation under rule labels than under labels assigned by
reading the turn.

The failure mode is intrinsic to trajectory mining: transcript formats are
undocumented, heterogeneous and evolving, and a parser that silently drops
content yields a corpus that looks fine, trains fine, and reports a stronger
effect than the truth. A class-level agreement breakdown between two independent
labelers is a cheap and unusually effective canary for it.

\section{Limitations and Discussion}

Human adjudication is single-annotator over 98 turns and not blind to the LLM's
label; we report no inter-human agreement, and the $\kappa = 0.727$ figure
assumes unadjudicated rule/LLM agreements are correct. A fresh 240-turn sample
from the final label set is released unadjudicated; adjudicating it is the
obvious next step. Paired tests over 5 folds have low power. The per-repository
cap samples turns, so a turn's immediate predecessor may not survive into the
corpus even though its history features were computed from the complete session
--- valid per row, but it disadvantages the prior-turn-text experiment. We did not
fine-tune a transformer, nor evaluate the LLM judge as a zero-shot
\emph{predictor} from the prompt alone --- an upper bound we would want before
concluding the task is hard rather than our models weak.

The workshop's premise is that we need better ways to describe what agents do.
Our experience suggests the descriptive step is not a preliminary to the
analysis; it dominates it. Our six categories are conventional, our two labelers
were both reasonable, and they still produced corpora supporting opposite
conclusions about how much of coding-agent work is debugging --- and rewriting two
sentences of the category definitions moved that conclusion again, with the
labeler held fixed. Reported behavior distributions in this literature are, we
suspect, less robust than they appear. We would go further than ``report an
agreement number'': a behavior taxonomy should be published together with the
exact definition text used to apply it, since that text is doing measurable work
that the category names conceal. On prediction: agent action type is weakly
anticipatable from developer language, mostly through session stickiness rather
than prompt semantics. We think that is worth reporting as a negative result, and
worth revisiting once labels exist that a predictor could be fairly scored
against.

\paragraph{Artifacts.} Corpus, three label sets, the human review sheet, all
result tables (including per-fold CV and paired significance tests), and the
pipeline: \url{https://github.com/medhab123/Prompt_Research}.

\begin{thebibliography}{9}
\small
\bibitem{tang2026programming}
Tang, Chen, Fang, Xu, Dhakal, Shi, McMillan, Huang, Li.
\emph{Programming by Chat: A Large-Scale Behavioral Analysis of 11,579
Real-World AI-Assisted IDE Sessions.} arXiv:2604.00436.

\bibitem{fang2026conversation}
Fang, Zhang, Tang, McMillan, Li, Huang.
\emph{From Conversation to Contribution: Characterizing Coding Agents in
Open-Source Software.} arXiv:2607.05677.

\bibitem{ong2024routellm}
Ong, Almahairi, Wu, Chiang, Wu, Gonzalez, Kadous, Stoica.
\emph{RouteLLM: Learning to Route LLMs with Preference Data.} arXiv:2406.18665.
\end{thebibliography}

\end{document}
